\begin{abstract}
We explore whether the Discrete Symplectic Cosmology (DSC) framework---in which cosmic evolution is modeled as a non-autonomous symplectic map on a Planck-scale lattice with a $1/\ln^2(t)$ adiabatic cooling law motivated by Riemann zero spectral statistics---can reproduce key statistical features of the cosmic microwave background (CMB) temperature anisotropies.
Using a St\"ormer--Verlet integrator on two- and three-dimensional lattices, we find that: (i) the symplectic evolution produces near-Gaussian one-point statistics (skewness $= +0.07 \pm 0.30$, excess kurtosis $= -0.28 \pm 0.24$ over 20 independent realizations), while a dissipative coupled-map-lattice baseline fails with kurtosis $= -1.60$; (ii) acoustic-like oscillation peaks emerge in the power spectrum $D(k) = k^2 P(k)$, with positions alignable to a toy analytic reference via three lattice parameters; (iii) the $1/\ln^2(t)$ cooling produces an effective acoustic freeze-out through asymptotic deceleration of wave propagation; and (iv) Bayesian optimization recovers all three lattice parameters from power spectra in a twin experiment (MSE $= 8.9 \times 10^{-5}$).
Full-sky maps generated by projecting a $128^3$ lattice onto a HEALPix sphere exhibit multi-scale temperature fluctuations with pixel distributions close to those of the Planck SMICA map, though quantitative comparison of the angular power spectrum yields $\chi^2/\mathrm{dof} \approx 3.3$, indicating that the forward model does not yet reproduce the observed $C_\ell$ at the level of individual multipoles.
A differentiable implementation in JAX enables gradient-based inverse reconstruction from Planck data (pixel correlation $r = 0.98$), but a held-out hemisphere test ($r = -0.17$) confirms that the severely underdetermined inverse problem does not uniquely constrain the 3D structure.
These results establish the DSC lattice as an exploratory proof-of-concept framework for generating CMB-like fluctuations from symplectic dynamics, while highlighting the substantial gap between phenomenological similarity and quantitative agreement with the observed CMB.
\end{abstract}
